% Options for packages loaded elsewhere
\PassOptionsToPackage{unicode}{hyperref}
\PassOptionsToPackage{hyphens}{url}
%
\documentclass[
  11pt,
]{article}
\usepackage{amsmath,amssymb}
\usepackage{iftex}
\ifPDFTeX
  \usepackage[T1]{fontenc}
  \usepackage[utf8]{inputenc}
  \usepackage{textcomp} % provide euro and other symbols
\else % if luatex or xetex
  \usepackage{unicode-math} % this also loads fontspec
  \defaultfontfeatures{Scale=MatchLowercase}
  \defaultfontfeatures[\rmfamily]{Ligatures=TeX,Scale=1}
\fi
\usepackage{lmodern}
\ifPDFTeX\else
  % xetex/luatex font selection
\fi
% Use upquote if available, for straight quotes in verbatim environments
\IfFileExists{upquote.sty}{\usepackage{upquote}}{}
\IfFileExists{microtype.sty}{% use microtype if available
  \usepackage[]{microtype}
  \UseMicrotypeSet[protrusion]{basicmath} % disable protrusion for tt fonts
}{}
\makeatletter
\@ifundefined{KOMAClassName}{% if non-KOMA class
  \IfFileExists{parskip.sty}{%
    \usepackage{parskip}
  }{% else
    \setlength{\parindent}{0pt}
    \setlength{\parskip}{6pt plus 2pt minus 1pt}}
}{% if KOMA class
  \KOMAoptions{parskip=half}}
\makeatother
\usepackage{xcolor}
\usepackage[margin=2.2cm]{geometry}
\setlength{\emergencystretch}{3em} % prevent overfull lines
\providecommand{\tightlist}{%
  \setlength{\itemsep}{0pt}\setlength{\parskip}{0pt}}
\setcounter{secnumdepth}{-\maxdimen} % remove section numbering
\ifLuaTeX
  \usepackage{selnolig}  % disable illegal ligatures
\fi
\IfFileExists{bookmark.sty}{\usepackage{bookmark}}{\usepackage{hyperref}}
\IfFileExists{xurl.sty}{\usepackage{xurl}}{} % add URL line breaks if available
\urlstyle{same}
\hypersetup{
  pdftitle={Worksheet 3A: Quadratics --- FOIL and Expanding},
  hidelinks,
  pdfcreator={LaTeX via pandoc}}

\title{Worksheet 3A: Quadratics --- FOIL and Expanding}
\author{}
\date{}

\begin{document}
\maketitle

\hypertarget{part-a-the-foil-method-basic-mechanics-12-problems}{%
\subsection{Part A: The FOIL Method --- Basic Mechanics (12
problems)}\label{part-a-the-foil-method-basic-mechanics-12-problems}}

\textbf{FOIL} stands for \textbf{F}irst, \textbf{O}uter, \textbf{I}nner,
\textbf{L}ast --- a method for multiplying two binomials.

\textbf{Example:} \((x + 3)(x + 2)\)

\begin{verbatim}
First:  x · x = x²
Outer:  x · 2 = 2x
Inner:  3 · x = 3x
Last:   3 · 2 = 6

Combine: x² + 2x + 3x + 6 = x² + 5x + 6
\end{verbatim}

\textbf{Expand each expression completely. Show your F-O-I-L steps for
the first few.}

\textbf{1.} \((t + 2)(t + 5)\)

\textbf{2.} \((p + 6)(p + 3)\)

\textbf{3.} \((m + 4)(m + 7)\)

\textbf{4.} \((r + 1)(r + 9)\)

\textbf{5.} \((k + 8)(k + 2)\)

\textbf{6.} \((n + 5)(n + 5)\) --- \emph{What do you notice about this
one?}

\textbf{7.} \((w + 3)(w + 3)\)

\textbf{8.} \((s + 10)(s + 1)\)

\textbf{9.} \((h + 6)(h + 6)\)

\textbf{10.} \((u + 4)(u + 4)\)

\textbf{11.} What pattern do you see when both binomials are identical,
like \((x + a)(x + a)\)?

\textbf{12.} Complete this rule: \((x + a)^2 = x^2 + \_\_\_\_\_ + a^2\)

\begin{center}\rule{0.5\linewidth}{0.5pt}\end{center}

\hypertarget{part-b-wait-exponents-dont-distribute-over-addition-8-problems}{%
\subsection{Part B: WAIT --- Exponents Don't Distribute Over Addition!
(8
problems)}\label{part-b-wait-exponents-dont-distribute-over-addition-8-problems}}

\begin{quote}
\textbf{Common Mistake Alert!}

Many students think \((a + b)^2 = a^2 + b^2\).

\textbf{This is WRONG.}

\((a + b)^2\) means \((a + b)(a + b)\), which you must FOIL out.
\end{quote}

\textbf{Let's prove this with actual numbers.}

\textbf{13.} Let \(a = 3\) and \(b = 4\).

\begin{itemize}
\item
  \begin{enumerate}
  \def\labelenumi{(\alph{enumi})}
  \tightlist
  \item
    Calculate \((a + b)^2 = (3 + 4)^2 = 7^2 ={}\) \_\_\_\_\_
  \end{enumerate}
\item
  \begin{enumerate}
  \def\labelenumi{(\alph{enumi})}
  \setcounter{enumi}{1}
  \tightlist
  \item
    Calculate \(a^2 + b^2 = 3^2 + 4^2 = 9 + 16 ={}\) \_\_\_\_\_
  \end{enumerate}
\item
  \begin{enumerate}
  \def\labelenumi{(\alph{enumi})}
  \setcounter{enumi}{2}
  \tightlist
  \item
    Are they equal? \_\_\_\_\_
  \end{enumerate}
\end{itemize}

\textbf{14.} Now let \(a = 2\) and \(b = 5\).

\begin{itemize}
\item
  \begin{enumerate}
  \def\labelenumi{(\alph{enumi})}
  \tightlist
  \item
    \((a + b)^2 ={}\) \_\_\_\_\_
  \end{enumerate}
\item
  \begin{enumerate}
  \def\labelenumi{(\alph{enumi})}
  \setcounter{enumi}{1}
  \tightlist
  \item
    \(a^2 + b^2 ={}\) \_\_\_\_\_
  \end{enumerate}
\item
  \begin{enumerate}
  \def\labelenumi{(\alph{enumi})}
  \setcounter{enumi}{2}
  \tightlist
  \item
    What's the difference between the two answers? \_\_\_\_\_
  \end{enumerate}
\end{itemize}

\textbf{15.} Expand \((a + b)^2\) using FOIL:

\((a + b)^2 = (a + b)(a + b) ={}\) \_\_\_\_\_\_\_\_\_\_\_\_\_\_\_\_\_

\textbf{16.} The correct expansion is \(a^2 + 2ab + b^2\). What's the
``missing piece'' when people incorrectly write \(a^2 + b^2\)?

\textbf{17.} Using \(a = 3\) and \(b = 4\) again: - (a) Calculate
\(2ab = 2(3)(4) ={}\) \_\_\_\_\_ - (b) Verify:
\(a^2 + 2ab + b^2 = 9 + 24 + 16 ={}\) \_\_\_\_\_ - (c) Does this match
\((a + b)^2 = 49\)? \_\_\_\_\_

\textbf{18.} Does the same mistake happen with subtraction? Test
\((a - b)^2\) vs \(a^2 - b^2\) with \(a = 5\), \(b = 2\): - (a)
\((a - b)^2 = (5 - 2)^2 ={}\) \_\_\_\_\_ - (b)
\(a^2 - b^2 = 25 - 4 ={}\) \_\_\_\_\_ - (c) Are they equal? \_\_\_\_\_ -
(d) What is the correct expansion of \((a - b)^2\)? \_\_\_\_\_

\textbf{19.} Here's another trap. Is \((a + b)^3 = a^3 + b^3\)? Test
with \(a = 2\), \(b = 1\): - (a) \((a + b)^3 = (2 + 1)^3 = 3^3 ={}\)
\_\_\_\_\_ - (b) \(a^3 + b^3 = 8 + 1 ={}\) \_\_\_\_\_ - (c) Are they
equal? \_\_\_\_\_

\textbf{20.} Complete the rule: \textbf{Exponents distribute over
\_\_\_\_\_\_\_\_\_\_\_\_\_ but NOT over \_\_\_\_\_\_\_\_\_\_\_\_\_.}

\emph{(Hint: \((ab)^2 = a^2b^2\) works, but
\((a + b)^2 \neq a^2 + b^2\))}

\begin{center}\rule{0.5\linewidth}{0.5pt}\end{center}

\hypertarget{part-c-foil-with-negative-numbers-12-problems}{%
\subsection{Part C: FOIL with Negative Numbers (12
problems)}\label{part-c-foil-with-negative-numbers-12-problems}}

When negatives appear, keep careful track of signs.

\textbf{Example:} \((x - 4)(x + 3)\)

\begin{verbatim}
First:  x · x = x²
Outer:  x · 3 = 3x
Inner:  -4 · x = -4x
Last:   -4 · 3 = -12

Combine: x² + 3x - 4x - 12 = x² - x - 12
\end{verbatim}

\textbf{Expand completely:}

\textbf{21.} \((t - 3)(t + 5)\)

\textbf{22.} \((p + 4)(p - 2)\)

\textbf{23.} \((m - 6)(m + 1)\)

\textbf{24.} \((r - 2)(r - 5)\)

\textbf{25.} \((k - 7)(k - 3)\)

\textbf{26.} \((n - 4)(n - 4)\) --- \emph{Squaring a binomial with
subtraction}

\textbf{27.} \((w - 8)(w - 8)\)

\textbf{28.} \((s + 5)(s - 5)\) --- \emph{What happens to the middle
term?}

\textbf{29.} \((h + 9)(h - 9)\)

\textbf{30.} \((u + 7)(u - 7)\)

\textbf{31.} What pattern do you see when the binomials are
\((x + a)(x - a)\)?

\textbf{32.} Complete this rule: \((x + a)(x - a) = \_\_\_\_\_\)

\begin{center}\rule{0.5\linewidth}{0.5pt}\end{center}

\hypertarget{part-d-foil-with-coefficients-10-problems}{%
\subsection{Part D: FOIL with Coefficients (10
problems)}\label{part-d-foil-with-coefficients-10-problems}}

When the variable has a coefficient, apply FOIL the same way.

\textbf{Example:} \((2x + 3)(x + 4)\)

\begin{verbatim}
First:  2x · x = 2x²
Outer:  2x · 4 = 8x
Inner:  3 · x = 3x
Last:   3 · 4 = 12

Combine: 2x² + 8x + 3x + 12 = 2x² + 11x + 12
\end{verbatim}

\textbf{Expand completely:}

\textbf{33.} \((2t + 1)(t + 3)\)

\textbf{34.} \((3p + 2)(p + 4)\)

\textbf{35.} \((2m - 3)(m + 5)\)

\textbf{36.} \((4r + 1)(r - 2)\)

\textbf{37.} \((3k - 2)(2k + 1)\)

\textbf{38.} \((5n + 2)(2n - 3)\)

\textbf{39.} \((2w + 3)(2w + 3)\) --- \emph{Squaring with coefficients}

\textbf{40.} \((3s - 1)(3s - 1)\)

\textbf{41.} \((4h + 5)(4h - 5)\) --- \emph{Difference of squares with
coefficients}

\textbf{42.} \((2u - 7)(2u + 7)\)

\begin{center}\rule{0.5\linewidth}{0.5pt}\end{center}

\hypertarget{part-e-patterns-worth-memorizing-8-problems}{%
\subsection{Part E: Patterns Worth Memorizing (8
problems)}\label{part-e-patterns-worth-memorizing-8-problems}}

These patterns appear so often that recognizing them instantly saves
time.

\hypertarget{pattern-1-perfect-square-trinomials}{%
\subsubsection{Pattern 1: Perfect Square
Trinomials}\label{pattern-1-perfect-square-trinomials}}

\[\boxed{(a + b)^2 = a^2 + 2ab + b^2}\]

\[\boxed{(a - b)^2 = a^2 - 2ab + b^2}\]

\emph{The middle term is always \textbf{twice} the product of the two
terms.}

\hypertarget{pattern-2-difference-of-squares}{%
\subsubsection{Pattern 2: Difference of
Squares}\label{pattern-2-difference-of-squares}}

\[\boxed{(a + b)(a - b) = a^2 - b^2}\]

\emph{The middle terms cancel out, leaving only the difference of
squares.}

\textbf{Use the patterns to expand quickly:}

\textbf{43.} \((t + 8)^2\)

\textbf{44.} \((p - 5)^2\)

\textbf{45.} \((m + 11)(m - 11)\)

\textbf{46.} \((3r + 4)^2\)

\textbf{47.} \((5k - 2)^2\)

\textbf{48.} \((6n + 1)(6n - 1)\)

\textbf{49.} \((2w + 9)^2\)

\textbf{50.} \((4s - 3)(4s + 3)\)

\begin{center}\rule{0.5\linewidth}{0.5pt}\end{center}

\hypertarget{part-f-discovering-pascals-triangle-12-problems}{%
\subsection{Part F: Discovering Pascal's Triangle (12
problems)}\label{part-f-discovering-pascals-triangle-12-problems}}

Let's expand powers of \((a + b)\) and look for patterns in the
coefficients.

\textbf{51.} \((a + b)^0 = 1\)

Coefficients: \textbf{1}

\textbf{52.} \((a + b)^1 = a + b\)

Coefficients: \textbf{1, 1}

\textbf{53.} Expand \((a + b)^2\) using FOIL.

\((a + b)^2 ={}\) \_\_\_\_\_\_\_\_\_\_\_\_\_

Coefficients of \(a^2\), \(ab\), and \(b^2\): \textbf{\_\_\_, \_\_\_,
\_\_\_}

\textbf{54.} Expand \((a + b)^3\) by computing
\((a + b) \cdot (a + b)^2\).

\emph{Use your answer from problem 53 and distribute \((a + b)\) through
it.}

\((a + b)^3 ={}\) \_\_\_\_\_\_\_\_\_\_\_\_\_\_\_\_\_\_\_

Coefficients: \textbf{\_\_\_, \_\_\_, \_\_\_, \_\_\_}

\textbf{55.} Expand \((a + b)^4\) by computing
\((a + b) \cdot (a + b)^3\).

\((a + b)^4 ={}\) \_\_\_\_\_\_\_\_\_\_\_\_\_\_\_\_\_\_\_\_\_\_\_\_

Coefficients: \textbf{\_\_\_, \_\_\_, \_\_\_, \_\_\_, \_\_\_}

\textbf{56.} Arrange your coefficients in a triangle:

\begin{verbatim}
                1                   (a + b)⁰
              1   1                 (a + b)¹
            1   2   1               (a + b)²
          1   3   3   1             (a + b)³
        1   ?   ?   ?   1           (a + b)⁴
\end{verbatim}

Fill in the missing numbers for row 4.

\textbf{57.} Look at how each number relates to the two numbers above
it. Describe the pattern:

\_\_\_\_\_\_\_\_\_\_\_\_\_\_\_\_\_\_\_\_\_\_\_\_\_\_\_\_\_\_\_\_\_\_\_\_\_\_\_\_\_

\textbf{58.} Using the pattern, predict row 5 (the coefficients for
\((a + b)^5\)):

\textbf{\_\_\_, \_\_\_, \_\_\_, \_\_\_, \_\_\_, \_\_\_}

\textbf{59.} This triangle is called \textbf{\_\_\_\_\_\_\_\_\_\_'s
Triangle}.

\emph{(Named after a 17th-century French mathematician, though
mathematicians in Persia, China, and India discovered it centuries
earlier.)}

\textbf{60.} Expand \((a + b)^5\) using your predicted coefficients:

\((a + b)^5 ={}\)
\_\_\_\_\_\_\_\_\_\_\_\_\_\_\_\_\_\_\_\_\_\_\_\_\_\_\_\_\_\_\_

\textbf{61.} \textbf{Bonus connection:} The number in row \(n\),
position \(k\) tells you ``how many ways can you choose \(k\) items from
\(n\) items?''

Row 4, position 2 has the number \textbf{6}.

Verify: List all ways to choose 2 letters from \{A, B, C, D\}:

\_\_\_\_\_\_\_\_\_\_\_\_\_\_\_\_\_\_\_\_\_\_\_\_\_\_\_\_\_\_\_\_\_\_\_\_\_\_\_\_\_

\emph{(There should be exactly 6 ways!)}

\textbf{62.} Using the triangle, how many ways can you choose 3 items
from 5 items?

\emph{(Look at row 5, position 3)}

\begin{center}\rule{0.5\linewidth}{0.5pt}\end{center}

\hypertarget{part-g-beyond-two-terms-8-problems}{%
\subsection{Part G: Beyond Two Terms (8
problems)}\label{part-g-beyond-two-terms-8-problems}}

For multiplying larger expressions, we distribute systematically.

\textbf{Example:} \((x + 2)(x^2 + 3x + 1)\)

\begin{verbatim}
Distribute x:   x·x² + x·3x + x·1    = x³ + 3x² + x
Distribute 2:   2·x² + 2·3x + 2·1    = 2x² + 6x + 2
Combine:        x³ + 3x² + 2x² + x + 6x + 2 = x³ + 5x² + 7x + 2
\end{verbatim}

\textbf{Expand completely:}

\textbf{63.} \((t + 3)(t^2 + 2t + 4)\)

\textbf{64.} \((p - 2)(p^2 + 4p + 1)\)

\textbf{65.} \((m + 1)(m^2 - m + 1)\) --- \emph{Something interesting
happens here!}

\textbf{66.} \((r - 1)(r^2 + r + 1)\) --- \emph{And here!}

\textbf{67.} \((2k + 1)(k^2 - 3k + 2)\)

\textbf{68.} \((n + 2)(n + 3)(n + 1)\) --- \emph{Do two at a time}

\textbf{69.} Expand \((w - 1)^3\) by first finding \((w - 1)^2\), then
multiplying by \((w - 1)\).

\textbf{70.} Expand \((2t + 1)^3\) using the same method.

\begin{center}\rule{0.5\linewidth}{0.5pt}\end{center}

\end{document}
